\section{相異なる八点の完全交叉から作る一基点境界}
\label{sec:reduced-quadric-boundary}
Boolean配置を固定しないため、三二次式の完全交叉そのものを点の評価で記述する。
この節ではKoszul障害の普遍性を仮定せず、フレームとその適用範囲を証明する。

\begin{lemma}[七点への二次評価と非零重み]\label{lem:reduced-octad-seven-evaluation}
三二次式の完全交叉$Z\subset\mathbb P^3$が八つの相異なる点$p_0,\ldots,p_7$からなるとする。
全点で非零な一次式$h$を取り、各代表を$h(p_i)=1$と正規化する。
次数2の評価像$F_2\subset k^8$は7次元であり、任意の七点への制限は同型である。
従って$F_2$を消す一次元の汎関数の重み$\lambda_0,\ldots,\lambda_7$は全て非零であり、
\[
 P\operatorname{diag}(\lambda_0,\ldots,\lambda_7)P^{\mathsf T}=0,
 \qquad P=(p_0\mid\cdots\mid p_7)
\]
が成立する。
\end{lemma}
\begin{proof}
$R=k[e_0,e_1,e_2,e_3]/(f_1,f_2,f_3)$とする。
完全交叉のHilbert関数は$(1,4,7,8,8,\ldots)$であり、$h$は$R$の非零因子である。
次数$n$の評価像$F_n$は$R_n$と同型で、$h=1$の正規化により$F_n\subset F_{n+1}$となる。
$F_3=k^8$であり、このフィルトレーションの随伴次数代数は$R/(h)$と同型である。
実際、$F_n/F_{n-1}$の核は$R_n$内で$hR_{n-1}$となる。
$R/(h)$は三変数の三二次式Artin完全交叉なので、Hilbert列は$(1,3,3,1)$、socleは次数3だけにある。

もしある一点のindicator $\delta_i$が$F_2$に入れば、最小の$d\leq2$について
$\delta_i\in F_d\setminus F_{d-1}$となる。
全ての$f\in F_1$に対して$\delta_if=f(p_i)\delta_i\in F_d$なので、
$\delta_i$の初期類は随伴次数代数の正次数元で消される。
これは次数$d\leq2$に非零socleを与え、矛盾する。
従って$\delta_i\notin F_2$である。
任意の七点へ制限した評価の核は、欠けた点のindicatorの倍だけなので零であり、次元7から同型となる。
$F_2$の唯一の消去汎関数が一点で重み零なら、そのindicatorが$F_2$に入るため、全重みは非零である。
最後の行列等式は全ての二次単項式への評価をまとめたものである。
\end{proof}

\begin{proposition}[一般の相異なる八点に対する一基点フレーム]
\label{prop:reduced-octad-one-point-frame}
前補題の完全交叉に二次式$q$を加え、$q(p_0)=0$、$q(p_s)\ne0$ $(s=1,\ldots,7)$とする。
$B=R/(q)$と置くと
\[
 (\dim B_n)_{n\geq0}=(1,4,6,4,1,1,1,\ldots).
\]
$P'=(p_1\mid\cdots\mid p_7)$の核の任意の三基底$k_0,k_1,k_2\in k^7$に対し、
\[
 b_j=\sum_{s=1}^7\frac{(k_j)_s}{q(p_s)}p_s^{\otimes3}\quad(j=0,1,2),
 \qquad b_3=p_0^{\otimes3}
\]
は$B_3^*$の基底である。
この基底から三因子対称テンソル$w$を作れば、$w\in Z_{\mathrm{sym}}$である。
さらに$w$が有限Frobenius条件を通るときは$w\notin S$である。
\end{proposition}
\begin{proof}
七点への一次評価も単射なので$\operatorname{rk}P'=4$であり、核は3次元である。
八点の三次評価テンソルは独立である。
上の三つの重みベクトルは$q$の値を掛けた後に$P'$の核になるので、$qR_1$を消す。
基点のcubeもこれを消す。
前三ベクトルは七非基点の評価上で独立であり、最後のcubeは基点成分だけを持つので、四つは独立である。
$q:R_1\to R_3$は単射であり、$\dim B_3=8-4=4$からこれらは基底をなす。
前補題により$q:R_2\to R_4$の階数は7なので$\dim B_4=1$である。
$n\geq5$では$R_{n-2}$の全八点評価が同型となり、$q$乗法の階数は7である。
低次数の値も合わせて所要のHilbert関数を得る。

閉包所属を示すため$q_t=q+t h^2$とし、前三基底では分母を$q(p_s)+t$へ変える。
$P'$は全射なので、ある$l_1,\ldots,l_7$について
$p_0+\sum_{s=1}^7l_sp_s=0$となる。
第3基底を
\[
 b_3(t)=p_0^{\otimes3}
       +t\sum_{s=1}^7\frac{l_s}{q(p_s)+t}p_s^{\otimes3}
\]
とすれば、全ての$q_tR_1$を消し、$t=0$で上の基底へ特殊化する。
一般の$t$では四二次式が基点自由完全交叉なので、対応するテンソルは$S$に入る。
よって$w\in Z_{\mathrm{sym}}$である。
有限Frobenius点での非対称化性は、因子2だけの可逆対称化が最初の二因子の関係空間を変えず、
基点を持つ$B$を完全対称Frobenius双対にしてしまうという、系\ref{cor:boolean-frobenius-not-symmetric}と同じ矛盾による。
\end{proof}

\begin{proposition}[八点の直交行列表示]\label{prop:reduced-octad-orthogonal-form}
上記の相異なる八点完全交叉は、点の順序、代表のスカラー、および共通射影基底を選べば
\[
 (I_4\mid M),\qquad MM^{\mathsf T}=I_4
\]
という列表示を持つ。
この表示は点配置の表示であり、indexed AS対象の完全な同型分類ではない。
また任意の直交行列から相異なる八点完全交叉が得られるとは主張しない。
\end{proposition}
\begin{proof}
四つの独立な点を最初の四列に選ぶ。
前補題の全て非零な重みについて、最初の四列を$\sqrt{\lambda_i}$倍、
後ろ四列を$\sqrt{-\lambda_i}$倍する。標数零の代数閉体なので選べる。
得た列行列を$(A\mid C)$とすれば$AA^{\mathsf T}-CC^{\mathsf T}=0$である。
$A$は可逆だから、左から$A^{-1}$を掛けると
$M=A^{-1}C$は$MM^{\mathsf T}=I_4$を満たす。
\end{proof}

\begin{remark}[古典的な自己付随八点との関係]
これはCayley octadの古典的なGale自己付随性に対応する表示である。
Plaumann--Sturmfels--Vinzant\cite[Theorem 3.10]{PSVQuartics}は
一般位置の八点について完全交叉と自己付随性の同値を述べている。
ここでは完全交叉から表示へ向かう方向を、評価フィルトレーションで証明した。
逆方向を退化した直交行列全体へ拡張せず、この古典的正規形そのものの新規性は主張しない。
\end{remark}

\begin{corollary}[最後の二次式の全係数域]\label{cor:octad-last-quadric-values}
前補題の八点を固定すれば、最初の三二次式を法とした第四二次式は
\[
 \{(z_0,\ldots,z_7)\in k^8:\ \sum_s\lambda_sz_s=0\}
\]
上の評価値と一意に対応する。
基点を$p_0$だけとする条件は$z_0=0$、$z_1\cdots z_7\ne0$であり、
非零スカラー倍を除けば五次元射影空間の開部分である。
これは二次式の係数域の次元であり、indexed同型類の次元を意味しない。
\end{corollary}
\begin{proof}
二次評価像は7次元の$\ker\lambda$に等しく、その核は最初の三二次式の空間である。
一点零条件はその中の一つの独立な線型条件なので、スカラー倍を除いた次元は5である。
残る七点での非消滅は開条件である。
\end{proof}

\begin{proposition}[残る対称化空間と四次双対]\label{prop:partial-symmetrizer-quartics}
$B=\operatorname{Sym}(V^*)/(L)$、$\dim B_3=4$とし、その三次双対基底から
因子$0,1,3$が対称なテンソル$w$を作る。このとき
\[
 S_2(w)\longrightarrow B_4^*,\qquad H\longmapsto H^{(2)}w
\]
は線型同型である。右辺は四階対称テンソルとして解釈する。
さらに$H$の階数は対応する四次テンソルの一因子縮約の階数に等しい。
\end{proposition}
\begin{proof}
$H^{(2)}w$は常に因子$0,1,3$について対称である。
$S_2$の方程式はさらに交換$(02)$についての対称性を要求し、合わせて完全対称性となる。
その三次縮約は全て$B_3^*$に属するので、完全対称な場合は$L\cdot\operatorname{Sym}^2(V^*)$を消し、
$B_4^*$に属する。逆に$F\in B_4^*$の各一次縮約は$B_3^*$に属する。
選んだ三次双対基底により各縮約を一意に表示すれば、$F=H^{(2)}w$となる行列$H$を一意に得る。
$w$の因子2から$B_3^*$への縮約写像は同型なので、最後の階数の等式も従う。
\end{proof}
\begin{corollary}[基点と階数1対称化写像]\label{cor:basepoint-rank-one-symmetrizer}
前命題の仮定の下で、$L$の共通零点$[p]\in\mathbb P(V)$と、
$S_2(w)$の非零階数1行列の射影類は一対一に対応する。
対応は$H^{(2)}w=p^{\otimes4}$で与えられる。
特に基点がないことは$S_2(w)$に非零階数1元がないことと同値である。
これは閉点の対応であり、非被約基点スキームの同定までは主張しない。
\end{corollary}
\begin{proof}
$[p]$が基点なら$p^{\otimes4}$は$L\cdot\operatorname{Sym}^2(V^*)$を消し、
$B_4^*$に属する。逆に非零対称四テンソルの一因子縮約の階数が1なら、
それを$p\otimes T$と書け、対称性から$T$も$p^{\otimes3}$のスカラー倍となる。
従ってそのテンソルは$c p^{\otimes4}$である。
これが$B_4^*$に属する条件は、全ての$q\in L$と二次式$r$について
$q(p)r(p)=0$となることであり、$p\ne0$なので全ての$q(p)=0$と同値である。
前命題の同型と階数の等式を適用すれば結論を得る。
\end{proof}
\begin{corollary}[三因子対称Frobenius候補の尽くされた二分岐]
\label{cor:partial-symmetric-exhaustive-split}
三因子対称な有限Frobenius候補を三次双対基底で表示する。
このとき次の二枝は互いに排反で、全候補を尽くす。
\begin{enumerate}
\item $S_2(w)$に非零階数1元がない。このとき$L$は基点を持たず、
四二次式完全交叉から来る完全対称化可能な枝に属する。
$S_2(w)$は可逆行列で生成される一次元空間である。
\item $S_2(w)$に非零階数1元がある。このとき$L$は基点を持ち、
$w$は完全対称化できない。
\end{enumerate}
第二枝のKoszul/AS正則性はこの二分岐だけからは決まらない。
\end{corollary}
\begin{proof}
前系により二条件は基点の有無そのものである。
基点がなければ、四変数多項式環の四二次式は高さ4の正則列となる。
その完全交叉の四次双対は一次元で、完全な乗法双対性により一因子縮約階数は4である。
三次双対基底をその四次生成元の縮約基底へ変えれば$w$は完全対称となる。
逆に基点を持つ有限Frobenius候補が完全対称化できれば、線型対称化判定から
$S_2$に可逆元があり、対応する四次テンソルが同じ$L$の完全交叉表示を与える。
これは基点に反する（定理\ref{thm:symmetric-frobenius-chart}）。
\end{proof}
\begin{proposition}[二つの対称化証明書を持つ全候補の分岐]
\label{prop:two-positive-global-split}
有限Frobenius候補$w$について、相異なる$j,\ell\in\{1,2,3\}$で
$S_j(w),S_\ell(w)$に可逆元があるとする。残る位置を$m$とすると、次のいずれか一方が成り立つ。
\begin{enumerate}
\item $w$は完全対称化可能であり、$S_m(w)$は可逆行列で生成される一次元空間である。
\item $S_m(w)$の全行列は特異であり、その中に非零階数1行列が存在する。
\end{enumerate}
この分岐は四因子独立の基底変更で不変である。
第二枝を排除するAS定理はここではまだ得られていない。
\end{proposition}
\begin{proof}
二つの証明書で三因子を対称化し、残る因子が位置2となるよう四窓を巡回移動する。
三つの対称因子のいずれを共通因子と見ても、未対称因子への交換方程式は同じ完全対称性を表す。
隣接縮約像の直交空間は$\bigwedge^2(V^*)+L$と書け、
対称な三因子での三重交差は$B_3^*$である。
有限Frobenius条件からその次元は4であり、一因子階数4によって$w$はこの空間の基底テンソルとなる。
従って系\ref{cor:partial-symmetric-exhaustive-split}を適用できる。
元の表示に戻る際には
$S_i((g_0\otimes\cdots\otimes g_3)w)=g_0 S_i(w)g_i^{-1}$
により、次元・行列階数・可逆元の有無は保たれる。
巡回移動で共通因子が変わるときも、三対称因子間の同一視で可逆な合成を施すだけなので、
残る対称化空間の階数条件は同じである。
\end{proof}

\begin{corollary}\label{cor:one-basepoint-symmetrizer}
一点基点の縮約族で$B_4^*=k p_0^{\otimes4}$となる場合、$S_2(w)$は一次元で、
その非零元は全て階数1である。従ってこのチャートの対称化失敗は行列式の消滅として直接検出できる。
\end{corollary}
\begin{proof}
前命題を適用する。$p_0^{\otimes4}$の一因子縮約の階数は1である。
\end{proof}

\begin{lemma}[一点基点の場合の低次socle]\label{lem:one-basepoint-cubic-socle}
$\dim B_3=4$かつ$B_4^*=k p^{\otimes4}$、$p\ne0$なら、
$B_3$内のsocleは三次評価$B_3\to k$の核であり、次元3である。
\end{lemma}
\begin{proof}
$B_4\simeq k$を$p^{\otimes4}$で同定すると、$\ell\in B_1$、$b\in B_3$の積は
$\ell(p)b(p)$に対応する。$\ell(p)\ne0$となる一次式があるので、
全ての一次式に殺されることと$b(p)=0$は同値である。
三次評価は非零であり、$B$は次数1で生成されるから、これがsocleそのものである。
\end{proof}

\begin{proposition}[三因子対称チャートでの四因子同値判定]
\label{prop:partial-symmetric-orbit-congruence}
四次元の二次式空間$L,L'$から三次双対基底で作る二つの三因子対称テンソル$w,w'$について、
$\dim B_3=\dim B'_3=4$、隣接二対二階数6、
\[
 S_1(w)=S_3(w)=kI,\qquad S_1(w')=S_3(w')=kI
\]
を仮定する。
このとき$w,w'$が四因子独立の$G$作用で同値であることと、$L,L'$が通常の$\operatorname{GL}(V^*)$同値であることは同値である。
従ってAS正則部分に対象を限れば、これは当該チャート内の添字保持同型の必要十分条件でもある。
一般の非正則候補については、ここでは四窓の表示同値を述べている。
\end{proposition}
\begin{proof}
$w'=(g_0\otimes g_1\otimes g_2\otimes g_3)w$とする。
線型symmetrizerの定義から
\[
 S_j(w')=g_0S_j(w)g_j^{-1}
\]
が従う。$j=1,3$の一次元条件により、$g_1,g_3$はそれぞれ$g_0$のスカラー倍である。
因子$0,1$の縮約像は、従って共通変換$g_0\otimes g_0$で移る。
その直交関係$\bigwedge^2(V^*)+L$を取れば$L,L'$の通常の同値を得る。

逆に共通の基底変換で$L$を$L'$へ移せば、三つの対称因子上の三次双対空間も移る。
双方は同じ四次元空間の基底なので、残る因子2の独立な可逆変換で基底を合わせられる。
従って$G$同値である。AS正則部分での最後の解釈は、標準四窓の同型判定による。
\end{proof}

\begin{proposition}[一般位置八点でも生じる非Koszul標本]
\label{prop:general-position-octad-nonkoszul}
全ての四点部分集合が独立である八点完全交叉からも、
有限Frobenius条件と全四次検査を通り、五次で七次元のKoszul障害を持つ境界が得られる。
これは当該標本の存在主張であり、一般点の障害次元の認定ではない。
\end{proposition}
\begin{proof}
反対称$4\times4$行列$A$の上三角成分を、順序$01,02,03,12,13,23$で
$(1,2,3,5,7,11)$とし、
\[
 M=(I-A)(I+A)^{-1},\qquad P=(I\mid M)
\]
とする。$MM^{\mathsf T}=I$であり、$P$の70個の四次小行列式は全て非零である。
八列の二次評価行列は階数7なので、その核は三つの二次式を与える。
$h=e_0+e_1+e_2+e_3$で非斉次化した三式の商の次元は8であり、
$h=0$上の斉次三式の商も8次元のArtin環となる。
従って無限遠に追加点はなく、三式の完全交叉は上の相異なる八点そのものである。
これらは有理数上のGroebner基底で検算できる。

列を$h=1$で正規化し、$q=e_1^2+2e_2^2+3e_3^2$とする。
これは$p_0=(1,0,0,0)$だけで零となる。
命題\ref{prop:reduced-octad-one-point-frame}から作ったテンソルについて、厳密消去は
\[
\begin{gathered}
\text{一因子階数 }(4,4,4,4),\quad\text{二対二階数 }(6,6,6),\\
(\dim W_{i,3},\dim W_{i,4})=(4,1)\quad\text{全始点},\qquad
\dim C_{i,i+4}=35\quad\text{全始点},\\
\operatorname{rk}d_2(0,5)=77,\quad\operatorname{rk}d_3(0,5)=36
\end{gathered}
\]
を与える。従って始点0の五次$H_2$は7次元である。
$p_0$の平方型関係による六次元cycle部分も単射であり、その商は一次元となる。
また$\dim S_1=\dim S_2=\dim S_3=1$、中央のsymmetrizerの階数は1である。
再現可能な入力、幾何の検査と同じ五次行列の検算は
\path{computations/verify_reduced_octad_frame.py}に置く。
\end{proof}
